\begin{abstract}
This paper describes our solution for the MaCVi USV Semantic Segmentation Challenge at CVPR 2026, which targets the MULTIAQUA dataset featuring RGB, thermal, and LiDAR modalities with a challenging nighttime test set.
We adapt the MemorySAM framework---which repurposes SAM2's temporal memory mechanism for cross-modal feature aggregation---to the maritime domain and investigate effective strategies for multi-modal fusion under severe illumination changes.
Specifically, we adopt soft mixture-of-experts gating in SAM2's LoRA layers to improve expert utilization when the number of experts is small, and introduce a quality-aware cross-modal fusion head that jointly compares all modality features to produce relative importance weights.
Unlike prior independent per-modality scoring that suffers from sigmoid saturation, our fusion head extracts quality-sensitive statistics---global average, global max, and channel-wise standard deviation---through modality-specific compression networks, enabling the gating to respond to per-image quality variation rather than converging to fixed weights.
To provide direct training supervision for the gating network, we attach a lightweight auxiliary segmentation head to each modality's backbone features and derive oracle gating weights from per-modality prediction quality; the fusion head is trained to match these oracle weights via KL divergence.
The resulting weights are used with max-normalization for memory modulation---preserving the best modality's features while suppressing degraded ones---and as raw softmax weights for output fusion.
To bridge the day-night domain gap without nighttime training data, we apply aggressive RGB corruption augmentations including extreme darkening, complementary random masking~\cite{shin2024crm}, and stochastic full-channel zeroing.
We empirically validate each design choice and report our results on the MULTIAQUA benchmark.
% Our method achieves XX.X\% mIoU on the challenge test set.
\end{abstract}
